\section{Conclusion}
\label{sec:conclusion}

Aspis now has a detailed Lean account of its private-spend relation, field and
circle arithmetic, released encoder order, four-round folding, distinct-query
sampling, transcript schedule, five-tree authentication, theft reduction, and
spent-marker state transition.  The formal probability work also puts false
acceptance in one finite experiment and proves an exact union bound with every
current failure branch listed once.

The strongest numerical result is conditional and work-normalized.  If the
protocol event is bounded by \(0.7\cdot2^{-100}\) and the eighteen external
events total at most \(0.3\cdot2^{-100}\), Lean proves false acceptance at
most \(2^{-100}\).  For an attacker who has already completed the work search,
the raw theft results retain a \(2^{-70}\) term plus explicit source,
cryptographic, credential, runtime, and victim-setup terms.  Those terms have
not been combined into a numerical deployed theft bound.

The next work is concrete: prove the remaining outer Rust equalities for the
transcript, five Merkle openings, relation/candidate consumer, spend state,
and close path; finish the exact event-level FRI and Fiat--Shamir application;
and state defensible primitive and runtime bounds.  A deployed privacy theorem
additionally needs the production masking, entropy, retry, commitment,
serialization, and observation model.

The archived program, proof, statement, build provenance, and finalized
transaction make the implementation being discussed identifiable.  The Lean
source makes much of its mathematical reasoning mechanically checkable.  The
remaining assumptions are stated in the theorem interfaces and in this
report, so completing the deployed claim is a list of exact proof obligations
rather than an appeal to unspecified implementation correctness.
